\subsection{代数的代数元素}

设 \(\mathbf{A}\) 是域 \(\mathbf{K}\) 上的一个代数，且 \(x\) 是 \(\mathbf{A}\) 的一个元素。有两种可能的情形：

\begin{enumerate}
\def\labelenumi{\alph{enumi})}
\item
  单项式族 \((x^{n})_{n\in\mathbb{N}}\) 在K上是自由的。这时我们称x在K上是超越的。存在多项式代数 K{[}X{]} 到 A 的由 x 生成的子代数 K{[}x{]} 上的一个同构，且后者在K上具有无限次数。
\item
  存在一个整数 \(n \geqslant 1\) ，使得这些单项式 \(1, x, \ldots, x^{n-1}, x^n\) 是线性相关的；这等同于说存在一个多项式 \(f \neq 0\) 属于 \(\mathbf{K}[\mathbf{X}]\) ，使得 \(f(x) = 0\) 。这时我们说 \(x\) 在 \(\mathbf{K}\) 上是代数的。满足上述性质的最小整数 \(n \geqslant 1\) 称为 \(x\) 在 \(\mathbf{K}\) 上的次数。如果 \(x\) 在 \(\mathbf{K}\) 上的次数是 \(n\) ，则单项式 \(1, x, \ldots, x^{n-1}\) 在 \(\mathbf{K}\) 上线性无关，并且存在元素 \(a_0, a_1, \ldots, a_{n-1}\) （属于 \(\mathbf{K}\)），使得
\end{enumerate}

\[
x ^ {n} = a _ {0} + a _ {1} x + \dots + a _ {n - 1} x ^ {n - 1}.
\]

多项式 \(f(\mathbf{X}) = \mathbf{X}^n - \sum_{k=0}^{n-1} a_k \mathbf{X}^k\) 是唯一的首一多项式，其次数为

\(n\) 的首一多项式，属于 \(\mathbf{K}[\mathbf{X}]\) ，使得 \(f(x) = 0\) ；它被称为 \(x\) 在 \(\mathbf{K}\) 上的极小多项式。

定理1. --- 设A是域K上的一个代数，x是A中在K上代数的元素，n是次数，f是x在K上的极小多项式。

\begin{enumerate}
\def\labelenumi{\alph{enumi})}
\item
  对于多项式 \(g \in \mathbf{K}[\mathbf{X}]\) ，要使它满足 \(g(x) = 0\) ，其充分必要条件是 \(g\) 是 \(f\) 的倍元。
\item
  映射 \(g \mapsto g(x)\) 通过过渡到商，定义了商代数 \(\mathbf{K}[\mathbf{X}] / (f)\) 到代数 \(\mathbf{K}[x]\) 上的一个同构，并且元素 1, x, \ldots, \(x^{n-1}\) 构成 \(\mathbf{K}[x]\) 在 \(\mathbf{K}\) 上的一组基。特别地， \([\mathbf{K}[x]: \mathbf{K}] = n\) 。
\item
  假设 A 是整环。则 K{[}x{]} 是域，且 f 是 K{[}X{]} 中唯一使得 \(f(x) = 0\) 的首一不可约多项式。
\item
  要使 \(x\) 在 \(\mathbf{A}\) 中可逆，其充分必要条件是 \(f(0) \neq 0\) ；于是我们有 \(x^{-1} \in \mathbf{K}[x]\) 。
\end{enumerate}

存在唯一的代数同态 \(\varphi: \mathbf{K}[\mathbf{X}] \to \mathbf{A}\) 使得 \(\varphi(\mathbf{X}) = x\) ；我们有 \(\varphi(\mathbf{P}) = \mathbf{P}(x)\) （对每个 \(\mathbf{P} \in \mathbf{K}[\mathbf{X}]\) ），并且 \(\varphi\) 的像等于 \(\mathbf{K}[x]\) 。设 \(\mathfrak{a}\) 是 \(\varphi\) 的核；根据构造，多项式 \(f\) （即 \(x\) 在 \(\mathbf{K}\) 上的极小多项式）属于 \(\mathfrak{a}\) ，并且它是 \(\mathfrak{a}\) 中次数最小的首一多项式。因此（IV，第～11页，命题11）我们有 \(\mathfrak{a} = (f)\) ，从而得到 \(a\) ）。断言 \(b\) ）立即由 \(a\) ）以及IV第～11页命题10的推论得出。

假设A是一个整环。则代数 K{[}x{]} 是一个整环且在K上具有有限次数，因此它是一个域（V，第～10页，推论）。于是 K{[}x{]} 的理想 (f) 是极大的，这意味着 f 在 K{[}X{]} 中不可约（IV，第～13页）。最后，设 g 是 K{[}X{]} 中一个使 \(g(x) = 0\) 的首一不可约多项式；由 a) 可知它是 f 的倍元，因此 g = f，这就证明了 c)。

还需证明 \(d\) ）。存在一个多项式 \(g \in \mathbf{K}[\mathbf{X}]\) ，其次数为 \(n - 1\) ，以及一个元素 \(a\) （属于 \(\mathbf{K}\)），使得 \(f(\mathbf{X}) = \mathbf{X}g(\mathbf{X}) + a\) ，从而 \(f(0) = a\) 。若 \(a = 0\) ，则我们有 \(xg(x) = f(x) = 0\) ，而 \(g(x) \neq 0\) ，于是 \(x\) 在 \(\mathbf{A}\) 中不可逆。另一方面，若 \(a \neq 0\) ，则我们有 \(x \cdot [-a^{-1}g(x)] = 1\) ，所以 \(x\) 在 \(\mathbf{A}\) 中可逆，且 \(x^{-1} = -a^{-1}g(x)\) 。

推论1. --- 设A是域K上的一个代数。对于A的元素x，它在K上是代数的当且仅当由x生成的A的子代数 \(K[x]\) 在K上为有限次。特别地，若A在K上为有限次，则A的每个元素都在K上代数。

推论2. —— 设E是K的一个扩域，A是E上的一个代数，x是A中在K上代数的元素。则x在E上是代数的，x在E上的极小多项式整除x在K上的极小多项式，且x在E上的次数至多等于x在K上的次数。

设 \(f\) 为 \(x\) 在 \(\mathbf{K}\) 上的极小多项式；我们有 \(f(x) = 0\) 且 \(f \in \operatorname{E}[\mathbf{X}]\) ，因此 \(x\) 在 \(\mathbf{E}\) 上是代数的，且 \(f\) 是 \(x\) 在 \(\mathbf{E}\) 上的极小多项式的倍式（定理1，a））。

注记. --- 设 E 是域 K 的一个扩域，且 x 是 E 中一个元素，它是某个首一不可约多项式 \(f \in K[X]\) 的根。则定理 1，c）表明 f 是 x 在 K 上的极小多项式。

例子。--- * 1) 在复数域 C 中，数 i 在素域 Q 上是次数为 2 的代数数；因为如果 \(f(\mathbf{X}) = \mathbf{X}^{2} + 1\) ，则 \(f(i) = 0\) ，且 \(x^{2} + 1 \neq 0\) 对所有 \(x \in Q\) 成立，因此 \(i \notin Q\) 。于是域 \(\mathbf{Q}(i)\) 是 Q 的 2 次扩张；它由所有形如 \(a + bi\) 的数组成，其中a、b是有理数。同样地，i在实数域R上是2次代数数，且C是R的2次扩张。

\begin{enumerate}
\def\labelenumi{\arabic{enumi})}
\setcounter{enumi}{1}
\tightlist
\item
  设K是一个域，F 是 K 上的一元有理函数域 \(\mathrm{K}(\mathrm{X})\) 。设E为 F 的子域 \(\mathrm{K}(\mathrm{X}^{3})\) ；我们有 \(\mathrm{F} = \mathrm{E}(\mathrm{X})\) 且X在E上是代数的，因为它是多项式 \(Y^{3} - X^{3}\) 的根，该多项式属于环 \(E[Y]\) ；该多项式在 \(E[Y]\) 中不可约，因为若不然，它至少会有一个一次因子，从而存在两个非零多项式 \(u(\mathrm{X})\) , \(v(\mathrm{X})\) （属于 \(K[X]\) ）使得 \((u(\mathrm{X}^{3}))^{3} = \mathrm{X}^{3}(v(\mathrm{X}^{3}))^{3}\) ，这是荒谬的，因为若 m 和 n 分别是 u 和 v 的次数，这将意味着 \(9m = 9n + 3\) ，或 \(3m = 3n + 1\) 。因此，域 F 在 E 上的次数为 3，且 F 中的每个元素都可以唯一地写成线性组合 \(f(\mathrm{X}^{3}) + \mathrm{X}g(\mathrm{X}^{3}) + \mathrm{X}^{2}h(\mathrm{X}^{3})\) ，其中 f、g、h 是三个属于 \(\mathrm{K}(\mathrm{X})\) 的有理分式。
\end{enumerate}

*3) 在实数域 \(\mathbf{R}\) 中，可以证明 \(^1\) 数 \(\pi\) 在素域 \(\mathbf{Q}\) 上是超越的。*

